%%%%%%%%%%%%%%%%%%%%%%%%%%%%%%%%%%%%%%%%%%%%%%%%%%%%%%%%%%%%%%%%%%%%%%
% Lớp 11 | Lý thuyết đồ thị | Đồ thị
% Dạng 1C2?1-3: Câu hỏi tổng hợp
% Số câu: 9  (sinh tự động bởi tach_cau.py; từ nay thầy sửa trực tiếp file này)
%%%%%%%%%%%%%%%%%%%%%%%%%%%%%%%%%%%%%%%%%%%%%%%%%%%%%%%%%%%%%%%%%%%%%%

%%%%%%%%%%%%%%%% TRẮC NGHIỆM 4 PHƯƠNG ÁN %%%%%%%%%%%%%%%%

% nguồn: Ok_TracNghiemLop11_Lan1
\begin{ex}%[1C2B1-3]%[ma:a6658abb]
	Trong các dãy số $\left(u_n\right)$ sau đây, dãy số  nào là dãy số tăng?
	 \choice
		{$u_n=\sin  n$}
		{$u_n=n(-1)^n$} 
		 {$u_n=\dfrac{1}{n}$}
		{\True $u_n=2^{n+1}$} 
	 \loigiai{
		Xét dãy số $\left(u_n\right)$ cho bởi $u_n=2^{n+1}$.\\ 
		Ta có $u_n>0$ và $\dfrac{u_{n+1}}{u_n}=2>1$ nên $\left(u_n\right)$ là dãy  số tăng.
	} 
\end{ex}


%%%%%%%%%%%%%%%% TỰ LUẬN %%%%%%%%%%%%%%%%

% nguồn: Ok_TuanLuan_Toan11_Lan1
\begin{bt}%[1C2B1-3]%[ma:1d2429d1]
	Xét tính  tăng, giảm và bị chặn của mỗi dãy số $\left(u_n\right)$ sau, biết số hạng tổng quát:
	\begin{multicols}{3}
		\begin{enumerate}[a)]
			\item  $u_n=\dfrac{n^2}{n+1}$;
			\item $u_n=\dfrac{2}{n^5}$; 
			 \item $u_n=(-1)^n.n^2$.
		\end{enumerate}
	\end{multicols}
	\loigiai{ 
		\begin{enumerate} [a)]
			\item  $u_n=\dfrac{n^2}{n+1}$;\\
			Xét  hiệu $u_{n+1}-u_n=\dfrac{(n+1)^2}{n+2}-\dfrac{n^2}{n+1}=\dfrac{(n^3+3n^2+3n+1)-(n^3+2n)}{(n+1)(n+2)}=\dfrac{3n^2+n+1}{(n+1)(n+2)}>0$.
			 Vậy $\left(u_n\right)$ là dãy số tăng.\\
			Vì $u_n=\dfrac{n^2}{n+1}>0$ với mọi $n \in \mathbb{N^*}$ nên dãy số  $\left(u_n\right)$ bị chặn dưới bởi $0$.
			\item  $u_n=\dfrac{2}{n^5}$;\\
			Ta có $u_n=\dfrac{2}{n^5}>0$ và $\dfrac{u_{n+1}}{u_n}=\dfrac{n^5}{(n+1)^5}<1$. 
			Vậy  $\left(u_n\right)$ là dãy số giảm.\\
			$1=u_1>u_n=\dfrac{2}{n^5}>0$ với mọi $n \in \mathbb{N^*}$ nên dãy  số $\left(u_n\right)$ bị chặn.
			\item $u_n=(-1)^n.n^2$.\\ 
			 Ta có $u_1=-1$; $u_2=4$; $u_3=-9$;\\
			Vậy $\left(u_n\right)$ là dãy số không tăng không giảm, không bị  chặn.
		\end{enumerate}
	 }
\end{bt}

% nguồn: Ok_TuanLuan_Toan11_Lan1
\begin{bt}%[1C2B1-3]%[ma:cdaa6423]
	Cho dãy số thực dương $(u_n)$. Chứng  minh rằng dãy số $(u_n)$ là dãy số tăng khi và chỉ khi $\dfrac{u_{n+1}}{u_n} > 1$ với mọi $n\in\mathbb{N}^*$.
	\loigiai{ 
		Do dãy số $(u_n)$ là dãy  số thực dương nên $u_n > 0$ với mọi $n\in\mathbb{N}^*$.\\
		Suy ra với mọi $n\in\mathbb{N}^*$, ta có $u_{n+1}  > u_n \Leftrightarrow \dfrac{u_{n+1}}{u_n} > 1$.\\
		Vậy dãy số $(u_n)$ là dãy số tăng khi và chỉ khi $\dfrac{u_{n+1}}{u_n} > 1$ với  mọi $n\in\mathbb{N}^*$.
	 }
\end{bt}

% nguồn: Ok_TuanLuan_Toan11_Lan1
\begin{bt}%[1C2B1-3]%[ma:2f7b5fab]
	Trong các dãy số $(u_n)$  được xác định như sau, dãy số nào bị chặn dưới, bị chặn trên, bị chặn?
	\begin{enumEX}{3}
		\item $u_n = n^2  + 2$;
		\item $u_n = -2n +  1$;
		\item $u_n =  \dfrac{1}{n^2+n}$.
	\end{enumEX}
	\loigiai{ 
		\begin{enumEX}{1}
			\item Với mọi $n\in\mathbb{N}^*$, ta có $u_n = n^2 + 2 \ge 1^2 + 2 > 2$.  \\
			Vậy dãy số $(u_n)$ là dãy số  bị chặn dưới.
			\item Với mọi $n\in\mathbb{N}^*$, ta có $u_n = -2n+1 \le  -2\cdot 1 + 1 < 0$. \\
			Vậy dãy số $(u_n)$ là dãy số bị chặn trên. 
			\item Với mọi $n\in\mathbb{N}^*$, ta có $n^2 + n \ge 1^2  + 1 > 1 > 0$ nên $0 < \dfrac{1}{n^2+n} < 1$. \\
			Vậy dãy số  $(u_n)$ là dãy số bị chặn.
		\end{enumEX}
	 }
\end{bt}

% nguồn: Ok_TuanLuan_Toan11_Lan1
\begin{bt}%[1C2K1-3]%[ma:23bd9bce]
	Xét tính tăng, giảm của mỗi dãy số $(u_n)$, biết 
	\begin{enumEX}{3}
		\item $u_n  = \dfrac{n-3}{n+2}$;
		\item $u_n  = \dfrac{3^n}{2^n\cdot n!}$;
		\item $u_n = (-1)^n \cdot (2^n +  1)$.
	\end{enumEX}
	\loigiai{ 
		\begin{enumEX}{1}
			\item  Ta có
			 \[
			u_{n+1} - u_n = \dfrac{n-2}{n+3} - \dfrac{n-3}{n+2} = \dfrac{5}{(n+2)(n+3)}  > 0,\ \forall n\in\mathbb{N}^*
			\] 
			hay  $u_{n+1} > u_n,\ \forall n\in\mathbb{N}^*$.
			Suy ra dãy số $(u_n)$ là dãy số  tăng.
			 \item Ta có
			\[ 
			u_{n+1} - u_n = \dfrac{3^{n+1}}{2^{n+1}\cdot (n+1)!} - \dfrac{3^n}{2^n\cdot n!} = \dfrac{3^n}{2^{n+1}\cdot (n+1)!} \left[3 - 2(n+1)  \right] = \dfrac{3^n(1-2n)}{2^{n+1}\cdot (n+1)!} <0,\ \forall n\in\mathbb{N}^*
			\] 
			 hay $u_{n+1} < u_n,\ \forall n\in\mathbb{N}^*$.
			 Suy ra dãy số $(u_n)$ là dãy số giảm.
			\item Ta  có dạng khai triển của dãy số $u_n$ là $-3$, $5$, $-9$, $17$, $\ldots$ nên dãy số $(u_n)$ không là dãy số tăng, không là dãy số giảm.
		\end{enumEX}
	 }
\end{bt}

% nguồn: Ok_TuanLuan_Toan11_Lan2
\begin{ex}%[1C2B1-3]%[ma:77a01e11]
	Xét tính  tăng, giảm và bị chặn của mỗi dãy số $\left(u_n\right)$ sau, biết số hạng tổng quát:
	\begin{multicols}{3}
		\begin{enumerate}[a)]
			\item $u_n=\dfrac{n^2}{n+1}$; 
			\item  $u_n=\dfrac{2}{n^5}$;
			\item $u_n=(-1)^n.n^2$. 
		\end{enumerate}
	\end{multicols}
	 \loigiai{
		\begin{enumerate} [a)]
			\item  $u_n=\dfrac{n^2}{n+1}$;\\
			Xét  hiệu $u_{n+1}-u_n=\dfrac{(n+1)^2}{n+2}-\dfrac{n^2}{n+1}=\dfrac{(n^3+3n^2+3n+1)-(n^3+2n)}{(n+1)(n+2)}=\dfrac{3n^2+n+1}{(n+1)(n+2)}>0$.
			Vậy $\left(u_n\right)$ là dãy số tăng.\\ 
			Vì $u_n=\dfrac{n^2}{n+1}>0$ với mọi $n \in \mathbb{N^*}$ nên dãy số $\left(u_n\right)$ bị chặn dưới bởi $0$. 
			 \item $u_n=\dfrac{2}{n^5}$;\\
			Ta có  $u_n=\dfrac{2}{n^5}>0$ và $\dfrac{u_{n+1}}{u_n}=\dfrac{n^5}{(n+1)^5}<1$.
			Vậy  $\left(u_n\right)$ là dãy số giảm.\\
			 $1=u_1>u_n=\dfrac{2}{n^5}>0$ với mọi $n \in \mathbb{N^*}$ nên dãy số $\left(u_n\right)$ bị chặn.
			\item $u_n=(-1)^n.n^2$.\\ 
			Ta có $u_1=-1$;  $u_2=4$; $u_3=-9$;\\
			Vậy $\left(u_n\right)$ là dãy số không tăng không  giảm, không bị chặn.
		\end{enumerate}
	} 
\end{ex}

% nguồn: Ok_TuanLuan_Toan11_Lan2
\begin{ex}%[1C2B1-3]%[ma:0a8b2bb1]
	Cho dãy số thực dương $(u_n)$.  Chứng minh rằng dãy số $(u_n)$ là dãy số tăng khi và chỉ khi $\dfrac{u_{n+1}}{u_n} > 1$ với mọi $n\in\mathbb{N}^*$.
	\loigiai{ 
		Do dãy số $(u_n)$ là dãy số thực dương  nên $u_n > 0$ với mọi $n\in\mathbb{N}^*$.\\
		Suy ra  với mọi $n\in\mathbb{N}^*$, ta có $u_{n+1} > u_n \Leftrightarrow \dfrac{u_{n+1}}{u_n} > 1$.\\
		Vậy dãy số $(u_n)$ là dãy số tăng khi và chỉ khi $\dfrac{u_{n+1}}{u_n} >  1$ với mọi $n\in\mathbb{N}^*$.
	} 
\end{ex}

% nguồn: Ok_TuanLuan_Toan11_Lan2
\begin{ex}%[1C2B1-3]%[ma:bd020d2b]
	 Trong các dãy số $(u_n)$ được xác định như sau, dãy số nào bị chặn dưới, bị chặn trên, bị chặn?
	\begin{enumEX}{3}
		 \item $u_n = n^2 + 2$;
		\item $u_n = -2n +  1$;
		\item  $u_n = \dfrac{1}{n^2+n}$.
	\end{enumEX}
	 \loigiai{
		\begin{enumEX}{1}
			\item Với mọi $n\in\mathbb{N}^*$, ta có $u_n = n^2 + 2 \ge 1^2 + 2  > 2$. \\
			Vậy  dãy số $(u_n)$ là dãy số bị chặn dưới.
			\item Với mọi $n\in\mathbb{N}^*$, ta  có $u_n = -2n+1 \le -2\cdot 1 + 1 < 0$. \\
			Vậy dãy số $(u_n)$ là dãy  số bị chặn trên.
			\item Với mọi $n\in\mathbb{N}^*$, ta có $n^2  + n \ge 1^2 + 1 > 1 > 0$ nên $0 < \dfrac{1}{n^2+n} < 1$. \\
			Vậy dãy số $(u_n)$ là dãy số bị chặn. 
		\end{enumEX}
	 }
\end{ex}

% nguồn: Ok_TuanLuan_Toan11_Lan2
\begin{ex}%[1C2K1-3]%[ma:faa534a3]
	Xét tính tăng, giảm của  mỗi dãy số $(u_n)$, biết
	\begin{enumEX}{3}
		\item $u_n  = \dfrac{n-3}{n+2}$;
		\item $u_n = \dfrac{3^n}{2^n\cdot n!}$; 
		\item  $u_n = (-1)^n \cdot (2^n + 1)$.
	\end{enumEX}
	 \loigiai{
		\begin{enumEX}{1}
			\item Ta có 
			\[ 
			u_{n+1}  - u_n = \dfrac{n-2}{n+3} - \dfrac{n-3}{n+2} = \dfrac{5}{(n+2)(n+3)} > 0,\ \forall n\in\mathbb{N}^*
			 \]
			hay $u_{n+1} > u_n,\ \forall n\in\mathbb{N}^*$. 
			Suy ra dãy số $(u_n)$ là  dãy số tăng.
			\item Ta có 
			\[ 
			u_{n+1} - u_n = \dfrac{3^{n+1}}{2^{n+1}\cdot (n+1)!} - \dfrac{3^n}{2^n\cdot n!} = \dfrac{3^n}{2^{n+1}\cdot (n+1)!} \left[3 - 2(n+1) \right] = \dfrac{3^n(1-2n)}{2^{n+1}\cdot (n+1)!} <0,\ \forall  n\in\mathbb{N}^*
			\] 
			hay $u_{n+1} < u_n,\  \forall n\in\mathbb{N}^*$.
			Suy  ra dãy số $(u_n)$ là dãy số giảm.
			\item Ta có dạng khai triển của dãy số $u_n$ là $-3$, $5$, $-9$, $17$,  $\ldots$ nên dãy số $(u_n)$ không là dãy số tăng, không là dãy số giảm.
		\end{enumEX}
	} 
\end{ex}
